\documentclass[10pt,a4paper]{amsart}
\usepackage[margin=19mm]{geometry}
\usepackage{amsmath,amssymb,mathtools}
\usepackage[scr=rsfs,cal=euler]{mathalfa}
\usepackage{xfrac}
\usepackage[unicode,hidelinks,bookmarks=false]{hyperref}
\newcommand{\ie}{i.e.\ }
\newcommand{\cf}{cf.\ }
\newcommand{\resp}{respectively\ }
\newcommand{\Subsection}{\subsection}
\providecommand{\nobreakdash}{\nobreak\hbox{-}}
\setlength{\emergencystretch}{2em}
\multlinegap0pt
\usepackage{bm}
\usepackage{bbm}
\usepackage{dsfont}
\usepackage{xspace}
\usepackage{cancel}
\usepackage{mathtools}
\usepackage{amsthm}
\makeatletter
\@ifundefined{lemma}{\newtheorem{lemma}{Lemma}}{}
\makeatother
\newcommand{\PP}{\mathbb{P}}
\title[The Increasing Gap Dynamics in a General Spatial Matching Mo]{The Increasing Gap Dynamics in a General Spatial Matching Model}
\author{Andr\'es Fielbaum, Roberto Cominetti, Jos\'e Correa}
\date{Source: arXiv:2602.09926v1 (CC BY 4.0)}
\begin{document}
\maketitle

\noindent\emph{Source and licence.} The Increasing Gap Dynamics in a General Spatial Matching Model. Version arXiv:2602.09926v1 (CC BY 4.0), distributed under the Creative Commons Attribution 4.0 International licence, as recorded in the arXiv licence metadata for this exact version and shown on the abstract page https://arxiv.org/abs/2602.09926v1. Full rights notice: licenses/arxiv-2602.09926-CC-BY-4.0.txt.\par\smallskip

\noindent\textit{Editorial scope.} Counted unit: the Lemma whose source label is \texttt{Lemma:TechnicalAppendix} in the article (source0.tex lines 991--994, source label \texttt{Lemma:TechnicalAppendix}), delivered together with the article's own complete original proof of it (source0.tex lines 996--1053, one \texttt{proof} environment of 58 lines). No mathematics is written, repaired or re-derived: statement and proof are the article's own text line for line. Required context retained uncounted: none. Every definition, display and label of those lines is retained unchanged. Omitted from this package and not redistributed: 1--990, 995--995, 1054--1184. Nothing inside a retained excerpt is omitted or altered: every retained body is delivered exactly as the article's lines are written, so no omission receipt of any kind accompanies this package. Citations: the retained excerpts carry no citation command at all, so no bibliography entry of the article is transcribed and no reference is redistributed. Reference adaptation: none was needed and none was made; every reference inside the delivered unit resolves inside it, so no unresolved reference or citation is produced. Printed slips kept literal: no printed word, symbol, hypothesis or number of the counted statement or of its proof was altered. Layout and class: the article's own class is \texttt{article} with packages including mathptmx, geometry, xcolor, array, bbding, multirow, amsmath, amssymb, amsthm, graphicx, natbib, lscape, url, hyperref; the wrapper is the lane's pinned \texttt{amsart} editorial wrapper at 10pt on A4, which restates the theorem-like environments the retained text needs and adds the article's own macro definitions that the retained text needs (\texttt{\textbackslash PP}), with extra wrapper packages (bm, bbm, dsfont, xspace, cancel, mathtools). The delivered numbering is the unit's own. Assets: no retained span contains a figure environment, a graphics-inclusion command, a PDF-page inclusion, a table or a TikZ picture; no image file is copied into this package, the package declares no asset dependency and no image file is redistributed with it. Licence: version arXiv:2602.09926v1 of this article is distributed under the Creative Commons Attribution 4.0 International licence, as recorded in the arXiv licence metadata for this exact version and shown on the abstract page https://arxiv.org/abs/2602.09926v1. A full rights notice is delivered at licenses/arxiv-2602.09926-CC-BY-4.0.txt. Characters: every retained excerpt is pure ASCII, so the delivered engine, XeLaTeX with its default Latin Modern text font, reports no missing character.
\begin{lemma}\label{Lemma:TechnicalAppendix}
    Let $C\in \beta(X^N)$ be any  Borel set with  positive Lebesgue measure $\mu_N(C)>0$. Then, there exists $\varepsilon>0$, which depends only on $C$, such that  $\PP(S^{2N+1}\in C | S^1=s) \geq \varepsilon$ uniformly for  $s \in X^N$.
    
\end{lemma}

\begin{proof}
We first argue that it suffices to consider the case where $C$ is a product $A_1 \times \ldots \times A_N$ of disjoint closed subsets $A_1,\ldots,A_N$ of $X$.
Indeed, by Lebesgue's density theorem almost every point in $C\setminus D$ has density 1 so that taking such a point $x \in C\setminus D$ and any $\kappa\in (0,1)$ we can find a small radius $r>0$ such that the $\|\cdot\|_\infty$-ball $R=\{s\in X^N:\|s-x\|_\infty\leq r\}$ satisfies 
$\mu_N(C \cap R) \geq \kappa\, \mu_N(R)$. Observe that $R$ is a product of rectangles $R=R_1 \times \ldots \times R_N$ and since $x\not\in D$ we can choose $r$ small enough so that the $R_k$'s are pairwise disjoint.
Setting $A_i=R_i\cap X$ it follows that the product set $\tilde C\triangleq A_1\times\ldots\times A_N$ 
is a Borel subset of $X^N\setminus D$, while
 the regularity of $X$ implies that $\mu_N(\tilde C)=\prod_{i=1}^N\mu(A_i)>0$.
Moreover, since $P_N$ and $Q_N$ have a strictly positive density it follows that  there exists some $\gamma>0$
such that
$\PP(S^{2N+1} \in C | S^1=s)\geq \gamma\; \PP(S^{2N+1} \in \tilde C | S^1=s)$. Hence, it suffices to find a uniform lower bound for the right hand side, that is,  for  sets  of product form $A_1\times\ldots\times A_N$
for pairwise disjoint $A_k$'s  with $\mu(A_k)>0$ and therefore $P(A_k)>0$ as well as $Q(A_k)>0$.

 Now, considering a potentially smaller radius $r>0$, we can further assume that there is a second family of pairwise disjoint closed subsets $B_1,\ldots,B_N$ included in the interior of $X$, which are disjoint from all the $A_k$'s and such that $\mu(B_k)>0$ (so that
$P(B_k)>0$ and $Q(B_k)>0$). Without loss of generality we may assume that the sets $B_k$ have a small diameter $\delta>0$ with $\delta$  smaller than the distance between any two of these sets $B_k$ and also smaller than the distance between these sets and  the boundary of $X$.
With this reduction we can assume from now on that $C$ is of the form $C=A_1 \times \ldots \times A_N$ with all these properties.
The rest of the proof consist in two steps: first, we show that starting from any initial state $s\in X^N$, after $N$ periods there is a uniformly positive probability that the chain reaches an intermediate state $S^{N+1}$ in which every $B_k$ contains exactly one driver. Then after $N$ additional periods we show that with a uniformly positive probability the chain reaches a state 
$S^{2N+1}\in A_1\times\ldots\times A_N$.

\noindent{\bf First step.} 
We claim that there exists $\varepsilon_1>0$, independent of the initial state $S^1=s$, such that with probability at least $\varepsilon_1$, after $N$ steps the drivers are relocated so that each $B_k$ contains exactly one driver.

Let $s\in X^N$ be an arbitrary initial state and let $\eta_k^1$ be the number of drivers in $B_k$. Denote $\eta^1_{N+1}=N-\sum_{k=1}^N\eta_k^1$ the drivers located in $X\setminus\cup_{k=1}^NB_k$. The superindex represents the time step $t=1$, as these values will evolve. We will show that there exists $\varepsilon_1'=\varepsilon_1'(\eta^1)$ that only depends on $\eta^1$ such that with probability at least
$\varepsilon_1'$ after $N$ steps
every set $B_k$ contains exactly one driver.
Since there is a finite number of possible configurations for $\eta^1$, by taking $\varepsilon_1$ as the minimum of these $\varepsilon_1'$ this property will hold uniformly for $s\in X^N$.  

In order to find  $\varepsilon_1'$ as required, we relocate the drivers one at a time  until we have $\eta^{N+1}_k=1$ for all $k=1,\ldots,N$ and $\eta^{N+1}_{N+1}=0$. In every iteration $t$, we have some sets $B_k$ that have exactly one driver. We call these the ``achieved'' sets with ``achieved'' drivers. When this happens for every $k=1,...,N$ we have 
completed our task. If this is not yet the case, we take any $j$ such that $\eta^t_j=0$. We want to remove a driver that is not  achieved, and replace it by a new driver located within $B_j$. We now prove that this occurs with a probability $\varepsilon_1''>0$ that 
only depends on the family of currently achieved sets $B_k$. Since there are finitely many combinations for such achieved sets, it then suffices to define $\varepsilon_1'$ as the product of all such $\varepsilon_1''$. 

Let $p=\min_{k}P(B_k)$ so that the probability that the new driver appears in any $B_j$ is at least $p$. Notice that this $p$ is strictly positive and does not depend on $s$ nor $\eta^1$. What about the probability of selecting a driver that is not achieved? It suffices to show that there exists $v_1>0$ such that the Lebesgue measure of the Voronoi regions covered by non-achieved drivers  is larger than $v_1$, since then the probability of the user being assigned to a non-achieved driver is at least $\rho\,v_1$ and we may then take $\varepsilon_1''=p\,\rho\,v_1$. In order to find an appropriate $v_1$ let us set $r=\delta/4$ and let $v$ be the $d$-volume of a half ball of radius $r$. We claim that we can take $v_1=\min\{\alpha(r),v\}$. In order to prove this, it suffices to 
 consider the case where there is a single non-achieved driver $s_o$ (if there are more, the joint volume of the Voronoi cells of non-achieved drivers would be larger). We distinguish two cases for $s_o$.
 \vspace{-2ex}
 
\begin{itemize}
    \item \textbf{Case 1:} If for every achieved $B_k$ we have $\mathop{\rm dist}(s_o,B_k)\geq\delta/2$, then  $B(s_o,r)\cap X$ is included in the Voronoi cell of $s_o$ which then has  a Lebesgue measure  at least $\alpha(r)\geq v_1$.
     \vspace{-1ex}

    \item \textbf{Case 2:} There is some achieved $B_k$ such that $\mathop{\rm dist}(s_o,B_k)<\delta/2$. Note that the choice of $\delta$ implies that this $k$ is unique and also that 
    $B(s_o,r)
    \subseteq X$.
    Let $s_k$ be the position of the only driver in $B_k$.  The Voronoi partition that considers just  $s_o$ and $s_k$ is given by a hyperplane passing through the mid point of the segment connecting $s_k$ and $s_o$. Now, at least half of the ball 
    $B(s_o,r)$
    is included in the half-space containing $s_o$, and since all the other achieved sets $B_i$ are at a distance larger than $\delta/2$ from $s_o$, this half ball is  included in the Voronoi cell of $s_o$ which then has a Lebesgue measure  of at least $v\geq v_1$.
\end{itemize}
 \vspace{-2ex}

\noindent Both cases combined show that regardless of the position of $s_o$, the described $\varepsilon''_1$ works. This concludes the first step: there is a lower bound $\varepsilon_1$, independent of $s$, such that $P(\eta^{N+1}_1=1,\ldots,\eta^{N+1}_N=1|S^1=s)\geq\varepsilon_1$.


\vspace{1ex}
\noindent{\bf Second step.} We now prove that the probability of moving from an arbitrary state $s$  having exactly one driver in each $B_k$ to a state $\tilde s\in A_1\times\ldots\times A_N$, also has a uniform lower bound $\varepsilon_2$ independent of $s$. 
Let $q>0$ be such that $Q(B_k)\geq q$ and $P(A_k)\geq q$ for all $k=1,\ldots,N$. Consider the driver at position $s_1$ and let  $B_{j}$ the set that contains $s_1$. Then, with probability $Q(B_{j})>q$ an  arriving user will appear in this set $B_j$ and will be assigned to the driver $s_1$, while with probability $P(A_1)\geq q$ the new driver will appear in a position $\tilde s_1\in A_1$.
Hence, with probability $q^2$ the state $s$ will transition to a new state with the first driver located in $A_1$. Proceeding sequentially by replacing the drivers at $s_2,s_3,\ldots$ with new drivers at $\tilde s_2\in A_2$, $\tilde s_3\in A_3, \ldots$
it follows that  after $N$ periods the state $s$ will transition to a state in $A_1\times\ldots\times A_N$, with probability at least $q^{2N}$.
\\[2ex]
Combining the first and second steps it follows that after $2N$ periods the state $S^{2N+1}$ has a probability at least $\varepsilon=\varepsilon_1 \varepsilon_2$ to belong to the target set $C=A_1\times\ldots\times A_N$, uniformly in the initial state $S^1=s$.
\end{proof}

\end{document}
